% ---------------------------------------------------------------------
\section{The Idea}
\label{sec:idea}
% ---------------------------------------------------------------------

\subsection{What HyperPod is}
HyperPod is a distributed LLM system that uses \textbf{consumer laptops as shared
computing resources}. Instead of renting a cloud GPU or needing one expensive
machine, several everyday laptops pool their compute and memory to run a model
together.

\subsection{The original architecture (the starting point)}
The project begins from an existing three-part design:

\begin{itemize}[leftmargin=1.4em,nosep]
  \item \textbf{Go} --- the master controller: HTTP API, task queue, worker
  heartbeats, and result collection. \Keep
  \item \textbf{Rust} --- the worker agents that connect to the controller and
  execute inference tasks. \Keep
  \item \textbf{llama.cpp} --- the LLM inference engine (originally prebuilt
  Windows binaries). \Keep, but used as a \emph{local} engine only (see
  \S\ref{sec:techstack}).
\end{itemize}

\subsection{The defining problem it solves}
The laptops belong to real people and sit on Wi-Fi, so they come and go: a lid
closes, Wi-Fi drops, a battery dies. HyperPod's defining capability is to
\textbf{survive that} --- to keep generating the answer even when a laptop
vanishes mid-sentence. This is the one thing the project is built around.

\subsection{What this guide deliberately is NOT about}
\begin{itemize}[leftmargin=1.4em,nosep]
  \item \Drop\ \textbf{Distributed fine-tuning / training.} It overlaps existing
  projects (Petals, QVAC Fabric), is far harder to make fault-tolerant, and adds
  nothing to the live demo. HyperPod is \textbf{inference only}. If personalization
  is wanted, a small LoRA adapter can be trained offline on one machine
  beforehand --- never distributed training during the event. \HYP
  \item \Drop\ \textbf{Beating anyone on raw speed.} Faster heterogeneous
  inference is already owned by prima.cpp (\S\ref{sec:landscape}). We do not
  compete there.
\end{itemize}
